\documentclass{article}
\input{../../../../lib.sty}

\title{\texttt{fn make\_int\_to\_bigint}}
\author{Michael Shoemate}
\date{}

\begin{document}

\maketitle

\contrib
Proves soundness of the implementation of \rustdoc{measurements/noise/nature/fn}{make\_int\_to\_bigint} in \asOfCommit{mod.rs}{f5bb719}.

\section{Hoare Triple}
\subsection*{Precondition}
\subsubsection*{Compiler-Verified}

\begin{itemize}
    \item Generic \texttt{T} implements trait \rustdoc{traits/trait}{Integer}
\end{itemize}

\subsubsection*{User-Verified}
None

\subsection*{Pseudocode}
\lstinputlisting[language=Python,firstline=2,escapechar=|]{./pseudocode/make_int_to_bigint.py}

\subsection*{Postcondition}
\begin{theorem}
    \validTransformation{(\texttt{T})}{\texttt{make\_int\_to\_bigint}}
\end{theorem}

\begin{proof}
    Since \texttt{IBig.from} is infallible, and subtraction by the constant
    \texttt{T::MIN\_FINITE} is also infallible, the function closure cannot raise
    data-dependent errors.

    The function always returns a vector of big integers of the same length as the input,
    so the output is always a member of the output domain.

    It remains to show stability.
    When the metric is non-modular, the function is simply the coordinatewise exact cast
    from \texttt{T} to \texttt{IBig}, so distances are unchanged.

    When the metric is modular, each coordinate is translated by the same constant
    $-\texttt{T::MIN\_FINITE}$.
    Let $m = 2^{\texttt{bits}(T)}$ be the modulus of \texttt{T}.
    The map $x \mapsto x - \texttt{T::MIN\_FINITE}$ is a bijection from the native integer
    representation of \texttt{T} onto $\{0, \ldots, m - 1\} \subset \mathbb{Z}$.
    For any two inputs $x, x'$ and any lift $z \in m\mathbb{Z}^d$,
    \[
        ((x - \texttt{T::MIN\_FINITE}) + z) - (x' - \texttt{T::MIN\_FINITE}) = (x + z) - x'.
    \]
    Therefore translating both arguments by the same constant does not change the set of
    admissible lifts in the definition of modular $L_p$ distance, so the modular distance
    is preserved exactly.

    Hence the transformation is 1-stable in either branch, satisfying the stability property.
\end{proof}


\end{document}
